\section{Experimental Setup}
\label{sec:experimental_setup}

We evaluate CodeGraph on SWE-bench Lite \cite{yang2024swebench}, a benchmark consisting of 300 real GitHub issues from 12 open-source Python repositories. The objective is repository-level retrieval: given only a natural-language issue description, the system must rank files by relevance, placing the files modified in the human patch at the top of the ranked list.

\subsection{Metrics and Protocol}

Although Personalized PageRank propagates relevance over fine-grained structural entities (classes, functions, and methods), we evaluate retrieval performance at the file level. File-level retrieval serves as the correct proxy for agent context acquisition, as coding agents typically load entire files into their context window to preserve syntactical context.

To map entity-level scores to file-level ranks, we group the ranked entities by their parent file and assign each file the score of its highest-ranked constituent entity. We measure performance using three metrics:
\begin{itemize}
    \item \textbf{Recall@$k$ (R@$k$):} The fraction of issues where the target file is retrieved within the top-$k$ results. R@10 is our primary metric, reflecting the typical context budget of AI coding agents.
    \item \textbf{Mean Reciprocal Rank (MRR):} The average of $1 / \operatorname{rank}$ for the target file, indicating the precision of the ranked results.
    \item \textbf{Zero-recall count:} The number of issues where the target file does not appear in the top-10 results, indicating a reachability failure.
\end{itemize}

\subsection{Baselines}

We evaluate CodeGraph against four baseline configurations to isolate the contributions of lexical matching, graph topology, and ranking propagation:
\begin{enumerate}
    \item \textbf{BM25-only (Basic):} A sparse lexical retriever executing over the entity signatures and docstrings. This baseline omits the dependency graph entirely, serving as a measure of standard text search.
    \item \textbf{One-hop ego graph:} A structural baseline that retrieves all nodes within a single hop of the BM25 seed nodes without ranking. This tests whether graph structure alone is sufficient without PageRank propagation.
    \item \textbf{CodeGraph (Initial):} An early system configuration that applied confidence-weighted PPR restart, representing the performance of graph-based ranking prior to our seed engineering interventions.
    \item \textbf{BM25-only (Seed-Engineered):} A lexical baseline that uses our optimized seed corpus (including compound tokenization and path hints) but omits PPR propagation, isolating the value of structural routing.
\end{enumerate}
